Da lei de atrito quadratica de Lanchester,
\[
\frac{dF}{dt}=-\lambda_EE,
\qquad
\frac{dE}{dt}=-\lambda_FF.
\]
Considere a combinacao
\[
L(t)=\lambda_FF(t)^2-\lambda_EE(t)^2.
\]
Derivando,
\[
\frac{dL}{dt}=2\lambda_FF\frac{dF}{dt}-2\lambda_EE\frac{dE}{dt}.
\]
Substituindo as equacoes diferenciais,
\[
\frac{dL}{dt}=2\lambda_FF(-\lambda_EE)-2\lambda_EE(-\lambda_FF)=0.
\]
Portanto, $L(t)$ e constante no tempo, isto e,
\[
\lambda_FF(t)^2-\lambda_EE(t)^2=\lambda_FF_0^2-\lambda_EE_0^2.
\]

Essa e exatamente a lei do quadrado de Lanchester na forma da eq. (5.58). Como a solucao do Problema 5.43 satisfaz o sistema diferencial, ela necessariamente preserva essa quantidade invariante.